\section{SP1}
\label{sec:sp1}

La Sección~\ref{sec:noir} concluyó que Noir permite describir la verificación como un programa,
pero que cada programa sigue dando lugar a un circuito específico. SP1 adopta el enfoque
alternativo presentado en la Sección~\ref{sec:computacion_prueba}, en el que el sistema de
restricciones describe un procesador y la prueba certifica que un programa arbitrario se ejecutó
correctamente sobre él. Para el desarrollador, el cambio consiste en escribir un programa
convencional en Rust y delegar en la máquina virtual la representación de su ejecución. Este
cambio no elimina la necesidad de especificar correctamente la relación que se pretende
demostrar, sino que la traslada al programa que se ejecuta. La presente sección describe el origen
y los componentes de SP1, la forma de escribir programas y de definir sus valores públicos, el
proceso que conduce desde el código fuente hasta una prueba verificable en Ethereum y las
limitaciones relevantes para un puente entre cadenas. La descripción corresponde a la serie~6 de
SP1, en particular a las versiones 6.1.0 y 6.2.2 que resuelven las dependencias del proyecto, ya
que esta serie introdujo cambios de arquitectura respecto de las anteriores.

% ---------------------------------------------------------------------------
\subsection{Origen y propósito}
\label{sec:sp1_origen}

SP1 es una máquina virtual de conocimiento cero de código abierto, desarrollada por Succinct
Labs, que permite demostrar la ejecución correcta de programas compilados para la arquitectura
RISC-V~\cite{sp1repo, sp1docsv6}. Su repositorio reconoce como antecedentes directos el conjunto
de herramientas criptográficas Plonky3, el diseño de tablas y de búsquedas cruzadas de la máquina
virtual Valida, la cadena de compilación de RISC Zero y la idea de los \textit{builtins} de Cairo,
que SP1 generaliza mediante precompilados implementados como tablas adicionales~\cite{sp1repo}.

La serie~6, publicada a partir de febrero de 2026, incorporó un sistema de pruebas denominado
Hypercube, basado en polinomios multilineales~\cite{sp1docsv6}. En esta serie, la máquina virtual
implementa el conjunto de instrucciones RV64IM, es decir, la base entera de 64 bits de RISC-V
junto con la extensión M de multiplicación y división~\cite{riscvmanual2024, sp1docsv6}, y los
programas se compilan para el destino \texttt{riscv64im-succinct-zkvm-elf}~\cite{sp1repo}. El
conjunto no incluye instrucciones de punto flotante, de modo que esas operaciones deben emularse
mediante software. Las versiones anteriores empleaban un conjunto de instrucciones de 32 bits, y
la documentación advierte que las cadenas de compilación de 32 y 64 bits son incompatibles entre
sí~\cite{sp1docsv6}. Los valores con los que opera el sistema de pruebas pertenecen al cuerpo
primo KoalaBear, de módulo $p = 2^{31} - 2^{24} + 1$, y a una extensión de grado~4 de dicho
cuerpo~\cite{sp1repo}.

Al igual que en el caso de Noir, conviene distinguir los componentes que intervienen. El programa
se escribe en Rust, aunque el repositorio indica que puede emplearse cualquier lenguaje compilado
mediante LLVM~\cite{sp1repo}. Su compilación al conjunto de instrucciones admitido se realiza con
una cadena de compilación de Rust distribuida por Succinct, que produce un ejecutable en formato
ELF. La máquina virtual ejecuta ese archivo y registra los datos necesarios para demostrar su
ejecución. El sistema de pruebas, formado por el protocolo Hypercube, un mecanismo de recursión y
una etapa final que traduce la prueba a Groth16 o PLONK, genera la prueba a partir de esos datos.
Por último, la verificación puede realizarse con el kit de desarrollo de SP1, escrito en Rust, o
mediante contratos verificadores en Solidity publicados por Succinct~\cite{sp1repo,
sp1contracts}.

La reutilización de código constituye el principal atractivo de este enfoque. El repositorio
indica que los programas pueden utilizar la biblioteca estándar de Rust y que la mayoría de las
bibliotecas de su ecosistema deberían ser compatibles~\cite{sp1repo}. Los programas del puente
utilizan, por ejemplo, bibliotecas existentes para SHA-256, Keccak-256, aritmética sobre la curva
BLS12-381 y codificación según el ABI de Solidity. Esta compatibilidad tiene, sin embargo,
límites. El programa se ejecuta sin sistema operativo, por lo que solo dispone de las llamadas al
sistema que define la propia máquina virtual, como la lectura de entradas, la publicación de
valores, la finalización de la ejecución y los precompilados descritos en la
Sección~\ref{sec:sp1_flujo}~\cite{sp1repo}. Las bibliotecas que dependen del acceso a archivos, a
la red o a otros servicios del sistema operativo no pueden utilizarse sin adaptación, y toda
dependencia debe poder compilarse para el destino de 64 bits mencionado. Que una biblioteca
compile, además, no implica que su costo de prueba resulte razonable, cuestión que se retoma en la
Sección~\ref{sec:sp1_limitaciones}.

SP1 debe distinguirse, asimismo, de la infraestructura de generación de pruebas que ofrece
Succinct. El kit de desarrollo permite generar las pruebas localmente, en procesador o en GPU, y
también delegarlas en una red de probadores remota, funcionalidad que requiere habilitar una
característica específica del kit y credenciales del servicio~\cite{sp1repo}. Los manifiestos del
proyecto no habilitan esa característica, por lo que su implementación genera las pruebas de forma
local.

El interés de SP1 para el presente trabajo reside en que materializa el enfoque basado en la
ejecución de programas que la hipótesis contrapone al de Noir. La lógica de verificación de un
puente, que incluye la decodificación de estructuras RLP, el cálculo de hashes y la verificación de
firmas, puede escribirse como un programa Rust convencional que reutiliza bibliotecas existentes,
mientras que la máquina virtual asume su representación como sistema de restricciones. Las
consecuencias de este enfoque sobre los tiempos de prueba no se anticipan en este capítulo, sino
que forman parte de la comparación del Capítulo~\ref{chp:zkbridge}.

% ---------------------------------------------------------------------------
\subsection{Programas, entradas y valores públicos}
\label{sec:sp1_programas}

Una aplicación de SP1 se compone de dos programas. El \textbf{programa invitado}
(\textit{guest}) es el que se compila para la máquina virtual y cuya ejecución se demuestra. El
\textbf{programa anfitrión} (\textit{host}) es un programa Rust convencional que se ejecuta fuera
de la máquina virtual y utiliza el kit de desarrollo para preparar las entradas, solicitar la
ejecución o la generación de la prueba y procesar el resultado~\cite{sp1repo, sp1docsv6}. La prueba
respalda únicamente la ejecución del programa invitado. Las operaciones del anfitrión, como la
obtención de datos de un nodo de la cadena de origen, su decodificación o la construcción de las
entradas, permanecen fuera de la ejecución demostrada, por lo que, si la corrección de esos datos
es relevante, el programa invitado debe comprobarla por sí mismo.

El anfitrión escribe las entradas en un objeto \texttt{SP1Stdin}, que las serializa, y el programa
invitado las lee mediante la función \lstinline|sp1_zkvm::io::read|, que las deserializa, o
mediante \lstinline|read_vec|, que devuelve los bytes sin procesar~\cite{sp1repo}. La
documentación del código advierte que este flujo de entrada está controlado por el
probador~\cite{sp1repo}. Las entradas cumplen, por lo tanto, el papel del testigo de la
Sección~\ref{sec:nocion_zk}, ya que nada garantiza que correspondan a datos reales salvo las
comprobaciones que el propio programa realice sobre ellas.

Las condiciones se comprueban con las construcciones ordinarias de Rust, como la macro
\lstinline|assert!| o el tratamiento de errores que interrumpe la ejecución. Una comprobación
fallida provoca la terminación del programa con un código de salida distinto de cero, y una
ejecución que termina de ese modo no puede dar lugar a una prueba que los verificadores acepten. La
función de verificación del kit comprueba por defecto que la ejecución haya terminado con éxito, y
el contrato verificador de la versión utilizada rechaza cualquier código de salida distinto de
cero~\cite{sp1repo, sp1contracts}.

Los valores públicos se publican mediante \lstinline|sp1_zkvm::io::commit|, que serializa un
valor, o mediante \lstinline|commit_slice|, que añade bytes sin procesar~\cite{sp1repo}. El
conjunto de bytes publicados forma parte del enunciado, se entrega al verificador junto con la
prueba y, en el contrato verificador, interviene a través de su hash SHA-256~\cite{sp1contracts}.
Su codificación debe acordarse con quien lo consume, de modo que un contrato en Solidity, por
ejemplo, necesita una codificación que pueda decodificar sin ambigüedad, como la del ABI de
Ethereum.

El programa cuya ejecución se demuestra se identifica mediante su clave de verificación, que el kit
deriva de forma determinista del ejecutable ELF durante la etapa de preparación~\cite{sp1repo}. En
el contrato verificador, esa clave se representa mediante un valor de 32 bytes que se proporciona
junto con los valores públicos y la prueba~\cite{sp1contracts}. Quien verifica debe fijar el valor
esperado, ya que, si no lo comprobara, aceptaría también una prueba válida de otro programa que
publicara valores con el mismo formato.

\begin{figure}[htbp]
	\centering
	\begin{minipage}{\textwidth}
		\begin{lstlisting}[language=Rust]
#![no_main]
sp1_zkvm::entrypoint!(main);

pub fn main() {
    // Entrada privada: la suministra el probador.
    let w = sp1_zkvm::io::read::<u64>();
    // Aritmetica comprobada: un desbordamiento detiene la ejecucion.
    let y = w
        .checked_mul(w)
        .and_then(|w2| w2.checked_mul(w))
        .and_then(|w3| w3.checked_add(w))
        .and_then(|s| s.checked_add(5))
        .expect("w fuera del dominio admitido");
    // Valor publico: pasa a formar parte del enunciado.
    sp1_zkvm::io::commit(&y);
}
\end{lstlisting}
	\end{minipage}
	\caption{Programa invitado de SP1 que calcula $y = w^{3} + w + 5$ con aritmética entera
	comprobada.}
	\label{fig:sp1_guest}
\end{figure}

Los datos que el programa lee y no publica no se entregan al verificador, pero esto no basta para
afirmar que permanezcan secretos. Por una parte, los valores públicos pueden permitir deducirlos.
Por otra, la propiedad de conocimiento cero depende del modo de prueba. Según el modelo de
seguridad de SP1, las pruebas Groth16 y PLONK, generadas con la biblioteca gnark, son de
conocimiento cero, mientras que las pruebas de núcleo y las comprimidas no satisfacen actualmente
esa propiedad~\cite{sp1securitymodel}. En concordancia con ello, el repositorio de la versión
utilizada incluye un componente experimental y no auditado destinado a añadir conocimiento cero a
los protocolos multilineales en los que se basan estas últimas~\cite{sp1repo}. En el diseño del
puente, la prueba que se verifica en la cadena de destino es una prueba Groth16, aunque la
propiedad de conocimiento cero no resulta necesaria en una aplicación cuyos datos son públicos
(Sección~\ref{sec:nocion_zk}).

La Figura~\ref{fig:sp1_guest} muestra un programa invitado que retoma la relación
$y = w^{3} + w + 5$ de los ejemplos de Circom y Noir. A diferencia de estos, que evalúan la relación
en $\mathbb{F}_p$, el programa emplea enteros de 64 bits sin signo, cuya aritmética no es la de un
cuerpo finito. Para evitar cualquier ambigüedad, cada operación se realiza de forma comprobada, de
modo que un resultado que exceda el rango de esos enteros interrumpe la ejecución en lugar de
reducirse módulo $2^{64}$. El programa acepta, por lo tanto, únicamente los valores de $w$ para los
que $w^{3} + w + 5 < 2^{64}$, es decir, $0 \leq w \leq 2\,642\,245$. En ese dominio, el resultado
coincide con el de los ejemplos anteriores, ya que el valor entero es menor que el módulo de
254~bits que estos emplean, pero fuera de él la relación no es equivalente, puesto que aquellos
aceptan cualquier elemento del cuerpo.

El programa lee $w$ como entrada privada, calcula $y$ con aritmética comprobada y publica $y$. La
Figura~\ref{fig:sp1_host} muestra el código anfitrión imprescindible para este flujo, que escribe
$w$ en la entrada, deriva las claves a partir del ejecutable, genera y verifica una prueba y lee el
valor público. La prueba respalda que el programa identificado por la clave de verificación,
ejecutado con alguna entrada, terminó correctamente y publicó $y$. El verificador recibe la clave
de verificación, los valores públicos y la prueba, pero no $w$. El ejemplo no ilustra, sin embargo,
ninguna forma de confidencialidad, ya que la función es estrictamente creciente en el dominio
admitido y el valor público $y$ determina, por lo tanto, un único $w$ que cualquiera puede calcular.

\begin{figure}[htbp]
	\centering
	\begin{minipage}{\textwidth}
		\begin{lstlisting}[language=Rust]
use sp1_sdk::{
    blocking::{ProveRequest, Prover, ProverClient},
    include_elf, Elf, ProvingKey, SP1Stdin,
};

const ELF: Elf = include_elf!("cubic-program");

fn main() {
    let client = ProverClient::from_env();
    let mut stdin = SP1Stdin::new();
    stdin.write(&3u64); // w: entrada privada

    let pk = client.setup(ELF).unwrap(); // claves derivadas del ELF
    let mut proof = client.prove(&pk, stdin).compressed().run().unwrap();
    client.verify(&proof, pk.verifying_key(), None).unwrap();

    let y: u64 = proof.public_values.read(); // valor publico
    println!("y = {y}");
}
\end{lstlisting}
	\end{minipage}
	\caption{Programa anfitrión que prepara la entrada, genera y verifica una prueba comprimida y
	lee el valor público.}
	\label{fig:sp1_host}
\end{figure}

En una compilación y ejecución realizadas para el presente trabajo con SP1~6.2.2, el ejecutable
del programa invitado resultó ser un archivo ELF para RISC-V de 64 bits. La ejecución con $w = 3$
publicó $y = 35$ tras 4802 instrucciones, cifra que incluye la lectura y deserialización de la
entrada, mientras que con $w = 3\,000\,000$ la comprobación de desbordamiento interrumpió el
programa, que terminó con código de salida~1. El programa anfitrión de la
Figura~\ref{fig:sp1_host} generó y verificó correctamente una prueba comprimida para $w = 3$.

% ---------------------------------------------------------------------------
\subsection{Compilación, ejecución, generación y verificación de la prueba}
\label{sec:sp1_flujo}

La Figura~\ref{fig:sp1_flujo} representa el recorrido desde el código fuente hasta una prueba
verificable. La primera etapa es la compilación del programa invitado, que el kit realiza con la
cadena de compilación de Succinct, ya sea desde el script de compilación del anfitrión o mediante
una orden específica, y que produce un ejecutable ELF para RISC-V de 64 bits~\cite{sp1repo}. El
anfitrión incorpora ese ejecutable en su propio binario. A partir de él, la etapa de preparación
deriva de forma determinista la clave de prueba y la clave de verificación, sin emplear ningún
valor secreto~\cite{sp1repo}. Esta preparación no debe confundirse con la configuración confiable
de algunos SNARKs descrita en la Sección~\ref{sec:snarks_starks}, ya que solo depende del programa.
Las entradas, por su parte, se preparan en el anfitrión y se codifican al escribirlas en el objeto
de entrada, con independencia de la preparación de las claves.

\begin{figure}[htbp]
	\centering
	\begin{tikzpicture}[
		herramienta/.style = {draw, rounded corners, align=center, minimum height=1.1cm,
			text width=2.2cm, font=\footnotesize},
		artefacto/.style = {draw, dashed, align=center, minimum height=0.9cm,
			text width=2.2cm, font=\footnotesize},
		flecha/.style = {-{Stealth[length=2.5mm]}, thick},
		dato/.style = {font=\footnotesize, align=center}
	]
		\node[artefacto] (fuente) at (0, 0) {Programa\\invitado (Rust)};
		\node[herramienta] (compilador) at (2.65, 0) {Compilación\\(cadena de\\Succinct)};
		\node[artefacto] (elf) at (5.3, 0) {Ejecutable\\ELF RISC-V};
		\node[herramienta] (setup) at (7.95, 2.5) {Preparación\\de claves};
		\node[artefacto] (claves) at (10.6, 2.5) {Claves de\\prueba y de\\verificación};
		\node[herramienta] (verify) at (13.25, 2.5) {Verificación\\(kit o contrato)};
		\node[herramienta] (exec) at (7.95, -2.2) {Ejecución\\(RV64IM)};
		\node[artefacto] (entradas) at (7.95, -4.0) {Entradas\\\texttt{SP1Stdin}};
		\node[herramienta] (core) at (10.6, -2.2) {Pruebas de\\fragmentos};
		\node[herramienta] (rec) at (13.25, -2.2) {Recursión y\\compresión};
		\node[herramienta] (wrap) at (13.25, 0.1) {Envoltura\\Groth16\\o PLONK};

		\draw[flecha] (fuente) -- (compilador);
		\draw[flecha] (compilador) -- (elf);
		\draw[flecha] (elf) -- (setup);
		\draw[flecha] (elf) -- (exec);
		\draw[flecha] (setup) -- (claves);
		\draw[flecha] (claves) -- (verify);
		\draw[flecha] (entradas) -- (exec);
		\draw[flecha] (exec) -- (core);
		\draw[flecha] (core) -- (rec);
		\draw[flecha] (rec) -- (wrap);
		\draw[flecha] (wrap) -- node[left, dato] {$\pi$ y valores\\públicos} (verify);
		\draw[flecha] (claves) -- node[left, dato] {clave de\\prueba} (core);
	\end{tikzpicture}
	\caption{Proceso de compilación, ejecución, generación y verificación de una prueba con SP1.}
	\label{fig:sp1_flujo}
\end{figure}

En la etapa de ejecución, la máquina virtual ejecuta el programa sobre las entradas. El kit
permite ejecutar sin generar ninguna prueba, en cuyo caso devuelve los valores públicos y un
informe con el número de instrucciones ejecutadas y de llamadas al sistema~\cite{sp1repo}. Cuando
se genera una prueba, la ejecución se divide además en \textbf{fragmentos} (\textit{shards}), de
aproximadamente $2^{22}$ instrucciones cada uno según la documentación~\cite{sp1docsv6}, y los
eventos registrados en cada fragmento dan lugar a la traza que se demuestra. La división acota los
recursos necesarios para probar cada fragmento y permite procesarlos de forma independiente.

La traza de cada fragmento se organiza en tablas, en las que cada tabla corresponde a un tipo de
instrucción o a un componente especializado y sus columnas contienen los operandos necesarios para
comprobar la operación~\cite{sp1docsv6}. Dado que una palabra de 64 bits no cabe en un elemento del
cuerpo KoalaBear, cada palabra se representa mediante cuatro valores de 16 bits~\cite{sp1repo}. Las
restricciones que relacionan las columnas de una misma tabla se comprueban mediante un protocolo de
tipo \textit{sumcheck}, denominado \textit{zerocheck}, y las \textbf{interacciones} entre tablas, por
ejemplo entre la tabla que registra una instrucción y las que verifican sus operandos, se comprueban
mediante LogUp-GKR, un argumento de derivadas logarítmicas que verifica que los valores enviados y
recibidos coinciden con sus multiplicidades~\cite{sp1docsv6, papini2023logupgkr}. La coherencia de
la memoria entre fragmentos se comprueba mediante un hash de multiconjuntos, que transforma cada
lectura y escritura en un punto de una curva elíptica mediante la función Poseidon2, de modo que la
igualdad entre los conjuntos de escrituras y de lecturas se reduce a la igualdad de sus
hashes~\cite{sp1docsv6}. Por último, el probador se compromete con toda la traza del fragmento como
un único polinomio multilineal, mediante un esquema de compromiso denominado \textit{jagged}, que
admite tablas de alturas distintas~\cite{hemo2025jagged}, construido sobre el esquema
BaseFold~\cite{zeilberger2023basefold, sp1repo}. Este esquema se basa en funciones de hash, por lo
que esta parte del sistema no requiere configuración confiable, y la interacción se elimina mediante
la transformación de Fiat-Shamir (Sección~\ref{sec:interactivas}).

Las pruebas de los fragmentos constituyen la prueba de núcleo, cuyo tamaño total es proporcional a
la extensión de la ejecución~\cite{sp1docsv6}. Para obtener una prueba de tamaño constante, el
mecanismo de recursión verifica las pruebas de los fragmentos dentro de la propia máquina de pruebas
y las combina en forma de árbol hasta obtener una única prueba \textit{comprimida}~\cite{sp1docsv6}.
Esta prueba puede verificarse con el kit, pero no está pensada para su verificación en Ethereum.
Para ello, una etapa adicional cambia la función de hash empleada por la recursión y, finalmente,
una etapa de envoltura produce una prueba Groth16 o PLONK sobre la curva BN254~\cite{sp1docsv6,
sp1repo}. Esta última etapa utiliza la biblioteca gnark~\cite{gnarkrepo} y puede ejecutarse mediante
una imagen de Docker o, habilitando una característica del kit, de forma nativa~\cite{sp1repo}.
Según la documentación, las claves del circuito de Groth16 proceden de la ceremonia Ignition,
complementada con contribuciones del equipo de Succinct, mientras que la variante PLONK reutiliza la
cadena de referencia universal de dicha ceremonia, y las pruebas resultantes ocupan unos 260 y
868 bytes, respectivamente~\cite{sp1docsv6}.

La verificación puede realizarse con el kit, que admite las pruebas de cualquiera de los modos y
comprueba por defecto que la ejecución terminó con éxito, o en Ethereum, con las pruebas Groth16 o
PLONK~\cite{sp1repo, sp1docsv6}. En este último caso, el contrato recibe la clave de verificación
del programa, los valores públicos y la prueba. Los primeros cuatro bytes de la prueba identifican la
versión del verificador, lo que permite a un contrato de enrutamiento dirigirla al verificador
correspondiente. El contrato de Groth16 de la versión de circuito 6.1.0, que corresponde a las
versiones utilizadas en el proyecto, comprueba la prueba con cinco entradas públicas, que son la
clave del programa, el hash SHA-256 de los valores públicos truncado a 253 bits para que quepa en el
cuerpo de BN254, el código de salida, que debe ser cero, una raíz que identifica las claves del
mecanismo de recursión y un valor de uso único~\cite{sp1contracts}.

La máquina virtual ofrece, además, \textbf{precompilados}, es decir, tablas especializadas que
demuestran operaciones costosas sin descomponerlas en instrucciones RISC-V. El programa los invoca
mediante llamadas al sistema y, en la versión utilizada, abarcan la permutación de Keccak, la
expansión y la compresión de SHA-256, operaciones sobre las curvas secp256k1, secp256r1, ed25519,
BN254 y BLS12-381, aritmética sobre los cuerpos de estas curvas y sobre enteros grandes, la función
Poseidon2 y la verificación de pruebas de SP1~\cite{sp1repo}. Su uso habitual no requiere escribir
estas llamadas, sino sustituir bibliotecas criptográficas comunes por versiones modificadas que
dirigen automáticamente sus operaciones hacia los precompilados~\cite{sp1docsv6}. Conviene
distinguir, por lo tanto, el precompilado, que forma parte del sistema de pruebas, de la biblioteca
modificada que lo invoca, y de las operaciones del anfitrión, que quedan fuera de la prueba. Si una
dependencia no utiliza la versión modificada, la operación se ejecuta como instrucciones ordinarias
y su costo aumenta considerablemente~\cite{sp1docsv6}.

Esta descripción permite identificar intervalos de tiempo de naturaleza distinta. La compilación
del programa invitado, la preparación de las claves y el despliegue del contrato verificador
constituyen una preparación que se realiza una vez por programa. La construcción del cliente del
kit y la obtención de los artefactos que necesita la etapa de envoltura son costos de puesta en
marcha. En cada operación se repiten, en cambio, la preparación de las entradas en el anfitrión,
que no forma parte de la prueba, la ejecución, la generación de las pruebas de los fragmentos, su
compresión y, cuando se requiere la verificación en Ethereum, la envoltura final, seguidas de la
verificación local y de la verificación en la cadena de destino, cuyo costo se expresa en gas y cuya
duración incluye la espera hasta la inclusión de la transacción. El modo de ejecución simulado del
kit, que produce valores públicos reales pero una prueba vacía, no genera ninguna prueba y no debe
confundirse con estas etapas. Ninguna de ellas tiene necesariamente un equivalente directo en Noir,
por lo que su correspondencia se establece en la metodología del Capítulo~\ref{chp:zkbridge}.

% ---------------------------------------------------------------------------
\subsection{Limitaciones y consideraciones de seguridad}
\label{sec:sp1_limitaciones}

La primera limitación se refiere, como en Noir, a la diferencia entre la expresividad del lenguaje
y la corrección de la especificación. La máquina virtual garantiza que el programa se ejecutó
correctamente, pero no que el programa compruebe lo que se pretende. Dado que las entradas están
controladas por el probador, un programa que reciba como entrada la raíz contra la que verifica la
inclusión de una transacción, en lugar de derivarla de un encabezado cuya autenticidad también
compruebe, demuestra únicamente la inclusión respecto de una raíz arbitraria. Las distinciones de
la Sección~\ref{sec:circom_limitaciones} entre inclusión, autenticidad y finalidad del bloque y
éxito de la operación siguen siendo responsabilidad del programa y de la aplicación que lo
utiliza.

La segunda consideración es la identificación del programa. Toda modificación del código, de sus
dependencias o de la cadena de compilación produce un ejecutable distinto y, en consecuencia, una
clave de verificación distinta, de modo que el verificador debe fijar la clave correspondiente a
una compilación concreta. Para que terceros puedan reproducir esa clave, la herramienta de
compilación ofrece la opción de compilar dentro de un contenedor de Docker con el fin de obtener
ejecutables reproducibles~\cite{sp1repo}. La verificación debe exigir, además, que la ejecución haya
terminado con éxito, condición que el kit y el contrato verificador comprueban por defecto.

SP1 admite también cálculo no restringido. El repositorio de la versión utilizada define una macro
que ejecuta un bloque de código fuera de las restricciones, descarta al finalizar cualquier cambio
en el estado de la máquina virtual y solo permite transmitir datos al programa a través de un flujo
de sugerencias que este lee después como una entrada más~\cite{sp1repo}. Su papel es análogo al de
las funciones no restringidas de Noir y al operador \lstinline|<--| de Circom, ya que los datos así
obtenidos solo quedan respaldados por las comprobaciones que el programa realice sobre ellos.

Los precompilados introducen supuestos propios. La documentación advierte que los precompilados de
curvas elípticas asumen que sus entradas son puntos válidos de la curva, sin comprobarlo, que los
valores de los cuerpos deben estar en su representación canónica y que ciertos casos, como el uso
del precompilado de suma para duplicar un punto o la presencia del punto en el infinito, deben
tratarse por separado para no producir pruebas imposibles de verificar~\cite{sp1docsv6}. Las
bibliotecas modificadas deben corresponder, además, a la versión de SP1 utilizada, y una
sustitución que no se aplique hace que la operación vuelva a ejecutarse como instrucciones
ordinarias sin que se produzca ningún error~\cite{sp1docsv6}.

El costo de la generalidad constituye la limitación central del enfoque. Cada instrucción ejecutada
ocupa filas en las tablas de la traza, cada palabra de 64 bits se representa mediante varios
elementos del cuerpo y las operaciones criptográficas que no utilizan precompilados se ejecutan como
secuencias de instrucciones ordinarias. Los hashes binarios y la aritmética sobre la curva
BLS12-381, presentes en la verificación de un puente, pueden dominar, en consecuencia, el número de
instrucciones de un programa. El costo final depende, además, de la versión, de la configuración de
los fragmentos, del hardware disponible, que incluye la posibilidad de emplear GPU, y del modo de
prueba solicitado, ya que la envoltura en Groth16 o PLONK añade una etapa a la generación. Por ese
motivo, cualquier medición debe asociarse a la versión, a la configuración y al hardware con que se
obtuvo.

Por último, la privacidad de las entradas depende del modo de prueba y de quién la genera. Como se
indicó en la Sección~\ref{sec:sp1_programas}, solo las pruebas Groth16 y PLONK son de conocimiento
cero, y esa propiedad protege las entradas frente al verificador, pero no frente a quien genera la
prueba, que necesariamente las conoce. Por ese motivo, la delegación de la prueba en un servicio
remoto implica, salvo en modalidades específicas que la documentación describe para ese fin, que el
servicio recibe las entradas del programa~\cite{sp1docsv6}. En un puente cuyos datos son públicos esta limitación carece de
consecuencias, pero debe tenerse en cuenta en otras aplicaciones.

Con SP1 concluye la exposición de los dos enfoques que el presente trabajo compara, el basado en
circuitos específicos, representado por Noir, y el basado en la ejecución de programas sobre una
máquina virtual, representado por SP1. La Sección~\ref{sec:otros_zk} reseña brevemente otras
herramientas del ecosistema, y el Capítulo~\ref{chp:zkbridge} describe la implementación del puente
con ambas tecnologías y la metodología con que se comparan.
